\section{Scaling Law y arquitectura del modelo}

\subsection{El nuevo cuello de botella del Scaling en el preentrenamiento}

\begin{warningbox}{Los datos de alta calidad se acercan a una constante}
La velocidad de crecimiento de los datos de entrenamiento de alta calidad es mucho más lenta que la del crecimiento del cómputo:
\begin{itemize}
    \item La cantidad de datos de texto de alta calidad puede considerarse aproximadamente una constante
    \item Los datos multimodales no logran mejorar de manera eficaz el «coeficiente intelectual» del texto
    \item Esto hace que Token Efficiency se convierta en una dirección de optimización más importante que Compute Efficiency
\end{itemize}
\end{warningbox}

\subsection{El optimizador Muon: rompe el dominio de 10 años de Adam}

El optimizador Adam dominó el entrenamiento de aprendizaje profundo durante casi 10 años. El optimizador Muon, al considerar las relaciones de dependencia de la estructura matricial entre los parámetros, logra una mejora de aproximadamente 2 veces en Token Efficiency:

\begin{itemize}
    \item Adam: considera cada elemento de los parámetros de manera independiente, ignorando las relaciones estructurales entre los parámetros
    \item Muon: considera la dependency entre los parámetros matriciales y optimiza mediante ortogonalización
    \item Efecto: el efecto de aprendizaje de 30T tokens de alta calidad equivale a que Adam aprenda 60T tokens
\end{itemize}

\subsection{Hallazgos sobre el Scaling en RL}

\begin{knowledgebox}{Diseño de Curriculum en el aprendizaje por refuerzo}
En el entrenamiento de Agent, no se puede hacer que el modelo resuelva de entrada la tarea más difícil (como demostrar una conjetura matemática no resuelta), porque de lo contrario la Sample Efficiency sería extremadamente baja. Una mejor manera es:
\begin{itemize}
    \item Entrenar con tareas de dificultad media
    \item Aumentar la dificultad de manera gradual (un proceso de ascenso de montaña)
    \item Combinarlo con una buena estrategia de Sampling
\end{itemize}
Esto es, en esencia, un proceso de Curriculum Learning.
\end{knowledgebox}

\subsection{Resumen de la sección}

El nuevo cuello de botella del Scaling en el preentrenamiento es la limitación de los datos de alta calidad. El optimizador Muon rompe el dominio de 10 años de Adam mediante la ortogonalización matricial, logrando una mejora de 2 veces en Token Efficiency. El entrenamiento de RL necesita un diseño de Curriculum para garantizar la Sample Efficiency.


